\documentclass[11pt]{article}
\usepackage[letterpaper,margin=0.85in]{geometry}
\usepackage[T1]{fontenc}
\usepackage{lmodern,microtype}
\usepackage{amsmath,amssymb,booktabs,array,xcolor,enumitem,fancyhdr,titlesec,hyperref}
\definecolor{navy}{HTML}{17324D}
\definecolor{teal}{HTML}{126C78}
\hypersetup{colorlinks=true,urlcolor=teal,linkcolor=teal,citecolor=teal,
 pdftitle={Sparse and Low-Rank Gaussian Graphical Models},pdfauthor={Xinyu Wang}}
\setlength{\parindent}{0pt}
\setlength{\parskip}{6pt}
\setlength{\headheight}{14pt}
\setlist[itemize]{leftmargin=1.3em,itemsep=3pt,topsep=3pt}
\titleformat{\section}{\large\bfseries\color{navy}}{\thesection}{0.6em}{}
\titleformat{\subsection}{\normalsize\bfseries\color{navy}}{\thesubsection}{0.6em}{}
\titlespacing*{\section}{0pt}{12pt}{5pt}
\titlespacing*{\subsection}{0pt}{8pt}{3pt}
\pagestyle{fancy}\fancyhf{}
\fancyhead[L]{\small\color{navy}Xinyu Wang\quad /\quad Research \& Technical Writing}
\fancyhead[R]{\small STAT 802}
\fancyfoot[L]{\footnotesize Literature review\;\textbullet\;Portfolio edition}
\fancyfoot[R]{\footnotesize\thepage}
\renewcommand{\headrulewidth}{0.3pt}
\newcommand{\tr}{\operatorname{tr}}
\newcommand{\rank}{\operatorname{rank}}
\newcommand{\supp}{\operatorname{supp}}
\newcommand{\sign}{\operatorname{sign}}
\begin{document}
\thispagestyle{fancy}
{\small\bfseries\color{teal}HIGH-DIMENSIONAL STATISTICS / CONVEX OPTIMIZATION}\par
\vspace{4pt}
{\LARGE\bfseries\color{navy}Sparse and Low-Rank\\[3pt]Gaussian Graphical Models}\par
{\large A technical review of latent-variable model selection}\par
\vspace{3pt}
\textbf{Xinyu (Ricky) Wang}\\
Statistics \& Data Science PhD candidate, University of Delaware\\
STAT 802 final report: May 2026\quad |\quad Portfolio revision: November 2026

\section*{Scope and contribution}
This report synthesizes the framework of Chandrasekaran, Parrilo, and Willsky
\cite{cpw}. My contribution is an exposition connecting the Schur complement,
sparse--low-rank decomposition, geometric identifiability, and statistical recovery.
The model, recovery results, and numerical studies reviewed here are attributed
to that paper. This edition does not present a new estimator, an original
consistency proof, or an independently executed experiment.

\section{Why latent variables change the graph}
A Gaussian graphical model describes conditional independence among jointly
Gaussian variables. For covariance $\Sigma\succ0$ and precision $K=\Sigma^{-1}$,
an off-diagonal zero $K_{ij}=0$ means that variables $i$ and $j$ are independent
conditional on the remaining variables. Precision entries are not themselves
correlation coefficients; the corresponding partial correlation is
$-K_{ij}/\sqrt{K_{ii}K_{jj}}$.

In high dimensions, the sample covariance can be singular or unstable. Sparse
precision estimation reduces complexity, but a sparse-only model can become
misleading when unobserved factors influence many observed variables. In stock
returns, for example, common market effects can induce widespread dependence.

Partition the joint precision matrix into observed ($O$) and hidden ($H$) blocks:
\[
K^*_{(OH)}=\begin{pmatrix}K_O^*&K_{O,H}^*\\K_{H,O}^*&K_H^*\end{pmatrix}.
\]
The Schur complement gives the \emph{marginal} precision of the observed vector:
\begin{equation}
\widetilde K_O^*=(\Sigma_O^*)^{-1}
=\underbrace{K_O^*}_{S^*}
-\underbrace{K_{O,H}^*(K_H^*)^{-1}K_{H,O}^*}_{L^*}.
\end{equation}
Here $S^*$ is the precision of $X_O$ \emph{conditional on} $X_H$. The matrix
$L^*\succeq0$ has rank at most the number of hidden variables and can be dense.
Thus, the observed precision $S^*-L^*$ need not be sparse even when the
conditional graph is sparse. Marginalization and conditioning are different
operations; this distinction is central to the interpretation of the model.

\newpage
\section{A joint estimation objective}
Let $\widehat\Sigma_O$ be the sample covariance of the observed variables.
The regularized likelihood formulation reviewed in \cite{cpw} is
\begin{equation}
\begin{aligned}
(\widehat S,\widehat L)=\arg\min_{S,L}\;&
-\log\det(S-L)+\tr\!\left[\widehat\Sigma_O(S-L)\right]\\
&+\lambda_n\bigl(\gamma\|S\|_1+\tr(L)\bigr),\\
\text{subject to }&S-L\succ0,\qquad L\succeq0.
\end{aligned}
\end{equation}
The likelihood term fits the observed covariance. The entrywise $\ell_1$ penalty
encourages sparsity in $S$, while the trace penalty encourages low rank in $L$.
For positive-semidefinite $L$, its trace equals its nuclear norm. The parameter
$\lambda_n$ controls overall regularization and $\gamma$ controls the relative
penalty on the sparse component.

The unconstrained regularizer
$f_\gamma(S,L)=\gamma\|S\|_1+\|L\|_*$ has dual norm
\begin{equation}
g_\gamma(S,L)=\max\left\{\frac{\|S\|_\infty}{\gamma},\|L\|_2\right\},
\end{equation}
where $\|\cdot\|_\infty$ denotes the largest absolute matrix entry and
$\|\cdot\|_2$ denotes the spectral norm. The placement of $\gamma$ matters:
the sparse-component norm is divided by $\gamma$.

\section{When can the two components be separated?}
Sparsity and low rank alone do not guarantee a unique decomposition. A matrix
concentrated on a small set of entries can be both sparse and low rank.
The geometric analysis asks whether local perturbations of the two structures
can mimic one another.

For a sparse matrix $S^*$ and a rank-$r$ matrix $L^*=UDV^\top$, the relevant
tangent spaces are
\begin{align}
\Omega(S^*)&=\{N:\supp(N)\subseteq\supp(S^*)\},\\
T(L^*)&=\{UY_1^\top+Y_2V^\top:Y_1,Y_2\in\mathbb R^{p\times r}\}.
\end{align}
The latter expression uses the ambient matrix space; symmetric perturbations
are used for the symmetric model. Local separability requires
$\Omega\cap T=\{0\}$, rather than orthogonality of the two spaces.

The quantities
\[
\mu(\Omega)=\max_{N\in\Omega,\,\|N\|_\infty\leq1}\|N\|_2,
\qquad
\xi(T)=\max_{N\in T,\,\|N\|_2\leq1}\|N\|_\infty
\]
measure how strongly perturbations can concentrate under the two norms.
One sufficient geometric criterion uses
\begin{equation}
\chi(\Omega,T,\gamma)=\max\left\{\frac{\xi(T)}{\gamma},\,2\mu(\Omega)\gamma\right\}.
\end{equation}
If $\mu(\Omega)\xi(T)<1/2$, choosing
$\xi(T)<\gamma<1/[2\mu(\Omega)]$ makes $\chi<1$.
This tangent-space condition is a useful starting point, but is not by itself
the full set of assumptions needed for statistical consistency.

\newpage
\section{From geometric separation to statistical recovery}
\subsection{Likelihood curvature and stability}
The likelihood adds a second layer of geometry. At the true marginal precision,
the Fisher information acts on a perturbation $\Delta$ as
\[
\mathcal I^*(\Delta)=\Sigma_O^*\Delta\Sigma_O^*.
\]
The recovery analysis must control both the information retained within the
structured tangent spaces and its leakage into complementary directions.
It also controls how the tangent space of the low-rank variety changes near
$L^*$. A very small nonzero singular value makes rank recovery less stable.
These are stronger requirements than merely separating two spaces at one point.

\subsection{A schematic reading of Theorem 4.1}
Write $\xi=\xi(T)$ and $\mu=\mu(\Omega)$. Let $\sigma$ be the smallest
nonzero singular value of $L^*$ and $\theta$ the smallest magnitude of a nonzero
entry of $S^*$. Under the identifiability and Fisher-information assumptions
of Proposition 3.3, with $\gamma$ in its admissible range, Theorem 4.1 of
\cite{cpw} uses the following scaling conditions:
\begin{align}
n&\gtrsim\frac{p}{\xi^4},
&\lambda_n&\asymp\frac{1}{\xi}\sqrt{\frac{p}{n}},\\
\sigma&\gtrsim\frac{\lambda_n}{\xi^2},
&\theta&\gtrsim\frac{\lambda_n}{\mu}.
\end{align}
This is a schematic summary, not a replacement for the full theorem.
The symbols $\gtrsim$ and $\asymp$ suppress factors involving the covariance
scale and the Fisher-information constants. Those factors, as well as the
admissible range of $\gamma$, must be retained when checking a particular model.

The theorem gives, with probability greater than $1-2e^{-p}$,
\begin{align}
\sign(\widehat S)&=\sign(S^*),
&\rank(\widehat L)&=\rank(L^*),\\
g_\gamma(\widehat S-S^*,\widehat L-L^*)&\lesssim
\frac{1}{\xi}\sqrt{\frac{p}{n}}.
\end{align}
The first line describes \emph{algebraic correctness}: recovering the signed
sparse structure and the rank of the latent component. The second controls
estimation error. Recovering a low-rank matrix does not identify a unique set
of named latent factors or a unique factorization.

\subsection{What the guarantee means in practice}
\begin{itemize}
\item \textbf{Structure and estimation are distinct targets.} A small matrix
error does not automatically imply exact edge or rank recovery when signals
are weak.
\item \textbf{Sample requirements depend on geometry.} The dimension $p$,
incoherence quantity $\xi$, and sparse structure all matter; a sample-size
condition depending only on $p$ omits essential assumptions.
\item \textbf{Theory guides rather than replaces validation.} Stability across
regularization choices and performance on held-out observations are separate
empirical questions.
\end{itemize}

\newpage
\section{Numerical evidence reported in the source paper}
\textbf{Provenance.} The studies and numerical values in this section are from
Section 6 of \cite{cpw}. They are reviewed here, not reproduced by a new run.

\subsection{Synthetic graphical models}
The source paper studies three 36-node settings: two cycle-based models with
two and three hidden variables, and a $6\times6$ grid with one hidden variable.
Each setting uses 50 trials. The reported recovery curves illustrate how
algebraic correctness improves as the sample-to-dimension ratio increases.
Such finite simulations illustrate the theory; they do not prove consistency.

\subsection{Stock-return example}
The empirical example uses 216 monthly observations from 1990--2007 for
84 S\&P 100 stocks. The paper compares a sparse-only fit with a latent-variable
fit. Its reported comparison is summarized below.
\begin{center}
\small
\begin{tabular}{@{}lrr@{}}
\toprule
\textbf{Reported quantity} & \textbf{Sparse-only} & \textbf{Sparse + latent}\\
\midrule
Sparse edges & 646 & 135\\
Latent components & 0 & 5\\
Parameter count (paper's convention) & 730 & 639\\
KL divergence relative to sample-covariance Gaussian & 44.4 & 17.7\\
\bottomrule
\end{tabular}
\end{center}
\textit{Source: Chandrasekaran et al. (2012), Section 6.2 and Figure 2.}

The parameter count $84+135+5\times84=639$ follows the paper's convention.
Lower divergence indicates better fit to its covariance-based reference,
not demonstrated out-of-sample forecasting or trading performance.

For orientation, if the reference Gaussian is $\mathcal N(0,\widehat\Sigma)$
and the fitted Gaussian is $\mathcal N(0,K^{-1})$, then
\[
D_{\mathrm{KL}}\!\left(\mathcal N(0,\widehat\Sigma)\,\|\,
\mathcal N(0,K^{-1})\right)
=\tfrac12\!\left[\tr(K\widehat\Sigma)-\log\det(K\widehat\Sigma)-p\right].
\]
This formula states the direction explicitly; KL divergence is asymmetric.

\section{Discussion and next steps}
The decomposition distinguishes local conditional relationships from shared
latent effects. Its use requires attention to Gaussian assumptions, structural
separation, and regularization stability. Temporal dependence and nonstationarity
also deserve investigation; graph edges alone do not establish causality.

A future implementation could compare support recovery, latent rank, estimation
error, and held-out likelihood against a sparse-only baseline. This remains
proposed follow-up work.

{\footnotesize\textbf{Figure note.} Original image attachments were unavailable;
see Figures 1--2 in the cited paper.}

\begin{thebibliography}{1}\small
\bibitem{cpw} V. Chandrasekaran, P. A. Parrilo, and A. S. Willsky (2012).
Latent variable graphical model selection via convex optimization.
\emph{The Annals of Statistics}, 40(4), 1935--1967.
\href{https://doi.org/10.1214/11-AOS949}{doi:10.1214/11-AOS949}.
\href{https://arxiv.org/html/1008.1290}{Open-access text}.
\end{thebibliography}
\end{document}
